\documentclass[12pt]{article}
% Préambule commun à tous les documents extraits de « La Revue CAPES »
\usepackage[T1]{fontenc}
\usepackage[utf8]{inputenc}
\usepackage{lmodern}
\usepackage{amsmath,amssymb,amsthm,amsfonts,mathrsfs}
\usepackage{setspace}
\usepackage{listings}
\usepackage{pgfplots}
\pgfplotsset{compat=1.18}
\usepackage{longtable}
\usepackage{adjustbox}
\usepackage{lettrine}
\usepackage{tkz-tab}
\usepackage{changepage}
\usepackage{pdflscape}
\usepackage{graphicx}
\usepackage{tikz}
\usepackage{xcolor}
\usepackage[a4paper,top=2cm,bottom=2.5cm,left=2.5cm,right=2cm]{geometry}
\usepackage{fancyhdr}
\usepackage{draftwatermark}
\SetWatermarkText{La RevueCapes}
\SetWatermarkLightness{0.9}
\SetWatermarkScale{2}
\usepackage{hyperref}
\hypersetup{colorlinks=true,linkcolor=blue!50!black,urlcolor=blue!50!black}

\definecolor{revue}{RGB}{20,60,120}

\pagestyle{fancy}
\fancyhf{}
\renewcommand{\headrulewidth}{0.4pt}
\setlength{\headheight}{15pt}

% \entete{Matière}{Titre}{Type de document}
\newcommand{\entete}[3]{%
  \fancyhead[L]{\small\textsc{La Revue CAPES} -- #1}%
  \fancyhead[R]{\small #2}%
  \fancyfoot[C]{\thepage}%
  \thispagestyle{plain}%
  \begin{center}
    {\color{revue}\rule{\textwidth}{1.2pt}}\\[4pt]
    {\small\textsc{Concours directs CAP-CEG / CAPES -- Burkina Faso}}\\[6pt]
    {\LARGE\bfseries\color{revue} #2}\\[6pt]
    {\large\itshape #1 -- #3}\\[2pt]
    {\color{revue}\rule{\textwidth}{1.2pt}}
  \end{center}
  \vspace{0.5cm}%
}

% Encadré pour les remarques ajoutées dans les corrigés
\newcommand{\remarque}[1]{\par\medskip\noindent\fcolorbox{revue}{revue!6}{\parbox{\dimexpr\linewidth-2\fboxsep-2\fboxrule}{\textbf{\color{revue}Remarque.} #1}}\par\medskip}


\begin{document}
\entete{Mathématiques}{Ancien sujet CAPES -- Session 2021}{Énoncé}
\begin{spacing}{1.5}

\textbf{Exercice 1 : (5pts)}

Étudier la nature( convergente ou divergente) des séries numériques suivantes : 

\[S_1=\sum^{+\infty}_{n=1}(1-\frac{1}{n})^n\] et \[S_2=\sum^{+\infty}_{n=1}\frac{n(n-1)}{(n+1)!}\].

\textbf{Exercice 2 : (5pts)}

1. Déterminer le développement en série entière de la fonction $u$ solution de l'équation différentielle : $(x-1)u^{\prime}(x) + u(x)=0$, avec $u(0)=-1$.

2. Déterminer le rayon de convergence de la série obtenue.

3. En déduire son domaine de convergence.

\textbf{Exercice 3 : (5pts)}

Calculer $\int\int_{D}f(x,y)dxdy$ dans les cas suivants :

1. $D= \lbrace y\geq 0,\ x+y\leq 1,\ y-x\leq 1\rbrace$ et $f(x,y)=x^{2}y$.

2. $D= \lbrace {0\leq y\leq x, x^2+y^2 \leq 1\rbrace}$ et $f(x,y)=e^{x^2+y^2}$.

\textbf{Exercice 4 : (5pts)}

Soit l'équation différentielle \\
$(E) : y^{\prime\prime}-4y^{\prime} + 4y = (t^2 - 1)e^t$.

1. Trouver une solution particulière de $(E)$ sous la forme $y_p= Q(t)e^t$ ou $Q$ est un polynôme de degré 2.

2. En déduire la solution générale de $(E)$.

3. Trouver la solution $h$ telle que $h(0)=h^{\prime}(0)=1$.

\end{spacing}
\end{document}
